
\documentclass{article}
\usepackage{colortbl}
\usepackage{makecell}
\usepackage{multirow}
\usepackage{supertabular}

\begin{document}

\newcounter{utterance}

\centering \large Interaction Transcript for game `cladder', experiment `full\_v1.5\_default', episode 5155 with qwen3.5{-}4b.
\vspace{24pt}

{ \footnotesize  \setcounter{utterance}{1}
\setlength{\tabcolsep}{0pt}
\begin{supertabular}{c@{$\;$}|p{.15\linewidth}@{}p{.15\linewidth}p{.15\linewidth}p{.15\linewidth}p{.15\linewidth}p{.15\linewidth}}
    \# & \multicolumn{2}{c}{Player} && \multicolumn{2}{c}{Game Master} \\
    \hline

    \theutterance \stepcounter{utterance}  
    & & & \multicolumn{4}{p{0.6\linewidth}}{
        \cellcolor[rgb]{0.9,0.9,0.9}{
            \makecell[{{p{\linewidth}}}]{
                \texttt{\tiny{[P1$\langle$GM]}}
                \texttt{You are an expert in causal inference. The following question is not a typical commonsense query, but rather a meticulously designed question created by a professor specializing in causal inference, intended to assess the students' mastery of the course content.} \\
\texttt{Imagine a self{-}contained, hypothetical world with only the following conditions, and without any unmentioned factors or causal relationships: Talent has a direct effect on elite institution admission status. Brown eyes has a direct effect on elite institution admission status. Method 1: We look at how talent correlates with brown eyes case by case according to elite institution admission status. Method 2: We look directly at how talent correlates with brown eyes in general. To understand how talent affects brown eyes, is it more correct to use the Method 1 than Method 2?} \\
\texttt{Start your answer with 'Yes' or 'No', followed by additional reasoning or evidence to support your explanation.} \\
            }
        }
    }
     \\ \\

    \theutterance \stepcounter{utterance}  
    & \multicolumn{4}{p{0.6\linewidth}}{
        \cellcolor[rgb]{0.561,0.737,0.561}{
            \makecell[{{p{\linewidth}}}]{
                \texttt{\tiny{[P1$\rangle$GM]}}
                \texttt{No} \\
\\ 
\texttt{**Reasoning:**} \\
\\ 
\texttt{In this hypothetical world, we are asked to understand the causal effect of **Talent** on **Brown Eyes**. However, the problem statement explicitly defines the causal structure as:} \\
\texttt{1.  Talent \$\textbackslash{}rightarrow\$ Elite Institution Admission} \\
\texttt{2.  Brown Eyes \$\textbackslash{}rightarrow\$ Elite Institution Admission} \\
\\ 
\texttt{There is **no direct causal link** stated between Talent and Brown Eyes. In causal inference, asking "how Talent affects Brown Eyes" when there is no direct arrow between them usually implies looking for an indirect effect (mediated through the admission status) or determining if they are independent.} \\
\\ 
\texttt{Let's analyze the two methods:} \\
\\ 
\texttt{*   **Method 2 (General Correlation):** This calculates the raw correlation between Talent and Brown Eyes across the entire population. Since both variables independently cause a third variable (Admission), they will likely be correlated due to **selection bias** (or collider bias) at the "Admission" node. If high Talent leads to admission and Brown Eyes leads to admission, the subset of people who get admitted will show a spurious correlation between Talent and Brown Eyes. This correlation reflects the structure of the selection process, not a causal influence of Talent on Brown Eyes. Since there is no causal path from Talent to Brown Eyes, the true causal effect is zero, but Method 2 will likely yield a non{-}zero number due to the confounding structure.} \\
\\ 
\texttt{*   **Method 1 (Case{-}by{-}case according to admission status):** This method looks at the relationship *within* the strata defined by the admission status (e.g., comparing those admitted vs. those not admitted, or looking at the conditional distribution).} \\
\texttt{    *   If we look at the population of **Admitted** students: Both Talent and Brown Eyes are positively selected. The correlation here is driven by the common effect (Admission), not causation.} \\
\texttt{    *   If we look at the population of **Non{-}Admitted** students: Both are negatively selected.} \\
\texttt{    *   Crucially, simply stratifying by the outcome (Admission) does not "block" the backdoor path or remove the selection bias; it often creates a **stratification bias** (Simpson's Paradox) if the relationship reverses between strata, or it simply highlights the spurious correlation caused by the collider.} \\
\\ 
\texttt{**The Core Causal Issue:**} \\
\texttt{To understand the causal effect of \$X\$ (Talent) on \$Y\$ (Brown Eyes) in a world where \$X \textbackslash{}rightarrow Z\$ and \$Y \textbackslash{}rightarrow Z\$, one must recognize that \$X\$ and \$Y\$ are likely independent or have no causal relationship. Neither Method 1 nor Method 2 directly measures a causal effect because there is no causal path.} \\
\\ 
\texttt{However, the question asks which is *more correct* to use to understand the effect.} \\
\texttt{*   **Method 2** gives a biased estimate of the association driven by the collider.} \\
\texttt{*   **Method 1** (stratifying by the collider) often exacerbates the bias or creates a misleading picture depending on how the strata are defined, but it does not magically reveal a causal link where none exists.} \\
\\ 
\texttt{Actually, the premise of the question contains a trick. If there is **no** causal relationship between Talent and Brown Eyes, the correct answer to "how does Talent affect Brown Eyes?" is **Zero**.} \\
\texttt{*   Method 2 will likely show a non{-}zero correlation (spurious).} \\
\texttt{*   Method 1 will show different correlations in different strata (Simpson's Paradox), potentially hiding the lack of a causal link or creating a false one.} \\
\\ 
\texttt{Therefore, relying on either method to infer a causal effect is flawed because the causal effect is null. However, if forced to choose which approach is theoretically sounder for *identifying* the lack of causation or avoiding the specific bias of the collider:} \\
\texttt{Method 2 is the standard observational association. Method 1 attempts to control for the outcome, which is generally incorrect practice in causal inference (collider stratification bias). You should not condition on a collider (Admission) to estimate relationships between its causes.} \\
\\ 
\texttt{Thus, neither method is "correct" for estimating a causal effect because the effect is zero. But if the question implies "Which method avoids the bias of the collider?", the answer is neither, as conditioning on the collider (Method 1) creates bias. If the question implies "Which method shows the raw association?", that is Method 2.} \\
\\ 
\texttt{**Re{-}evaluating the specific phrasing:** "To understand how talent affects brown eyes..."} \\
\texttt{Since there is no causal path, the effect is 0.} \\
\texttt{Method 2 (General correlation) will likely show a non{-}zero correlation due to the common effect (Admission).} \\
\texttt{Method 1 (Stratified) will show varying correlations.} \\
\\ 
\texttt{In standard causal inference courses, the lesson regarding \$X \textbackslash{}rightarrow Z\$ and \$Y \textbackslash{}rightarrow Z\$ is that \$X\$ and \$Y\$ are independent. Any correlation observed is spurious.} \\
\texttt{*   Method 2 captures the spurious correlation.} \\
\texttt{*   Method 1 (conditioning on \$Z\$) creates **collider stratification bias**, which can induce a correlation even if \$X\$ and \$Y\$ are independent, or reverse an existing one.} \\
\\ 
\texttt{Therefore, **neither** method correctly identifies the causal effect (which is zero) without further intervention or knowledge of the underlying probabilities. However, usually, in such "trick" questions, the intended answer highlights that **Method 2** is the naive association, while **Method 1** is the biased adjustment.} \\
\\ 
\texttt{Wait, let's look closer at the prompt's constraints: "without any unmentioned factors".} \\
\texttt{If \$T \textbackslash{}rightarrow A\$ and \$B \textbackslash{}rightarrow A\$.} \\
\texttt{Is there a causal effect \$T \textbackslash{}rightarrow B\$? No.} \\
\texttt{Does Method 1 help? No, conditioning on \$A\$ induces dependence between \$T\$ and \$B\$ if they were independent, or distorts it if they weren't.} \\
\texttt{Does Method 2 help? It shows the marginal association.} \\
\\ 
\texttt{Perhaps the question is a trick about the definition of "affect". If there is no causal link, the answer is "No, neither method works, or Method 2 is the only one that reflects the lack of direct influence (by showing a correlation that is explained by the third variable, whereas Method 1 creates artificial dependencies)."} \\
\\ 
\texttt{Actually, the most rigorous answer in a causal inference context is: **No**.} \\
\texttt{Reasoning: To estimate the causal effect of \$X\$ on \$Y\$, one must block backdoor paths. Here, there is no backdoor path between \$X\$ and \$Y\$ (they are independent causes of \$Z\$). The causal effect is zero.} \\
\texttt{*   **Method 2** measures the marginal association \$P(B|T)\$. This will be non{-}zero due to the common effect \$Z\$. It is a biased estimator of the causal effect (which is 0).} \\
\texttt{*   **Method 1** conditions on the collider \$Z\$. Conditioning on a collider opens a non{-}causal path between \$T\$ and \$B\$, creating a spurious correlation (or distorting it). This is a classic example of **collider stratification bias**.} \\
\\ 
\texttt{Therefore, using Method 1 is **incorrect** because it introduces bias by conditioning on the outcome. Using Method 2 is also incorrect for estimating the causal effect (since it yields a spurious correlation), but it is the standard observational measure. However, the question asks "is it more correct to use Method 1". Since Method 1 actively creates bias (collider stratification) where Method 2 simply reflects the marginal distribution (which is also biased for causal inference but doesn't *add* bias via conditioning), Method 1 is definitely **not** more correct. In fact, Method 1 is often *worse* than Method 2 for understanding the relationship between causes of a common effect.} \\
\\ 
\texttt{**Conclusion:** No. Method 1 involves conditioning on a collider (elite institution admission), which induces a spurious correlation between Talent and Brown Eyes (collider stratification bias). Since there is no direct causal link between Talent and Brown Eyes, the true causal effect is zero. Method 1 will likely suggest a non{-}zero effect (or a different one than the marginal), making it incorrect. Method 2 shows the marginal association, which is also not the causal effect, but Method 1 is specifically flawed by the act of stratification on the outcome. Thus, it is not more correct to use Method 1.} \\
\\ 
\texttt{**Final Answer Formulation:**} \\
\texttt{Start with "No".} \\
\texttt{Explain that there is no causal path from Talent to Brown Eyes.} \\
\texttt{Explain that Method 1 conditions on a collider (Admission), inducing spurious correlation (Collider Stratification Bias).} \\
\texttt{Explain that Method 2 reflects the marginal association driven by the common effect, but neither correctly identifies the causal effect (which is zero), yet Method 1 is actively biased by the stratification.} \\
\\ 
\texttt{Wait, is there a scenario where Method 1 is better? Only if we wanted to see the effect *conditional* on admission, but the question asks "how talent affects brown eyes" (general causal effect). Conditioning on the outcome is never the way to find a causal effect between causes of that outcome.} \\
\\ 
\texttt{Therefore, the answer is **No**.} \\
\\ 
\texttt{**Step{-}by{-}step derivation for the final output:**} \\
\texttt{1.  **Identify the Causal Graph**: \$Talent \textbackslash{}rightarrow Admission\$ and \$Brown Eyes \textbackslash{}rightarrow Admission\$.} \\
\texttt{2.  **Identify the Target**: Causal effect of \$Talent\$ on \$Brown Eyes\$.} \\
\texttt{3.  **Analyze the Path**: There is no directed path from \$Talent\$ to \$Brown Eyes\$. The causal effect is 0.} \\
\texttt{4.  **Analyze Method 2 (Marginal Correlation)**: This measures \$Cov(T, B)\$. Since both cause \$Admission\$, \$T\$ and \$B\$ are correlated via the collider \$Admission\$. This is a spurious correlation. It does not represent the causal effect.} \\
\texttt{5.  **Analyze Method 1 (Stratified Correlation)**: This measures \$Cov(T, B | Admission=1)\$ and \$Cov(T, B | Admission=0)\$. Conditioning on a collider (\$Admission\$) opens a non{-}causal path between \$T\$ and \$B\$. This is known as **Collider Stratification Bias**. It can create a correlation where none exists (if \$T\$ and \$B\$ were independent) or distort the existing spurious correlation.} \\
\texttt{6.  **Compare**: Method 1 is theoretically flawed for causal inference in this context because it conditions on a collider. Method 2 is the standard observational association. Neither correctly estimates the causal effect (0), but Method 1 is *more* incorrect because it actively generates bias through stratification. Therefore, it is **not** more correct to use Method 1.} \\
\\ 
\texttt{**Answer**: No.} \\
            }
        }
    }
    & & \\ \\

    \theutterance \stepcounter{utterance}  
    & & & \multicolumn{2}{p{0.3\linewidth}}{
        \cellcolor[rgb]{0.9,0.9,0.9}{
            \makecell[{{p{\linewidth}}}]{
                \texttt{\tiny{[GM$|$GM]}}
                \texttt{No} \\
            }
        }
    }
    & & \\ \\

    \theutterance \stepcounter{utterance}  
    & & & \multicolumn{2}{p{0.3\linewidth}}{
        \cellcolor[rgb]{0.9,0.9,0.9}{
            \makecell[{{p{\linewidth}}}]{
                \texttt{\tiny{[GM$|$GM]}}
                \texttt{no} \\
            }
        }
    }
    & & \\ \\

    \theutterance \stepcounter{utterance}  
    & & & \multicolumn{2}{p{0.3\linewidth}}{
        \cellcolor[rgb]{0.9,0.9,0.9}{
            \makecell[{{p{\linewidth}}}]{
                \texttt{\tiny{[GM$|$GM]}}
                \texttt{game\_result = WIN} \\
            }
        }
    }
    & & \\ \\

\end{supertabular}
}

\end{document}
